% §7 The characterization theorem — draft §7 (lines 232–292), PLUS Lemma 9.1 lifted to
% its tail (draft lines 326–338; NOTES-exposition Part 4 repair 2: "k is the memory
% kernel" belongs where the reader meets k, and §8's Gamma/stable examples use it).
% Add the forward reference from Remark 7.6 to hypothesis (H) — NOTES Part 1 item 7:
% the extreme boundary of the cone IS the boundary of (H); §9 meets the endpoint three
% independent times.
% Blueprint part 07; shared labels: lem:selfdecomposable-increment (and -derivative),
% thm:main-characterization (with thm:main-construction / thm:main-analysis /
% prop:main-uniqueness as its halves), cor:semigroup-case, rem:drift-boundary,
% rem:cone-geometry (check exact label), lem:memory-kernel.
% NOTE (CLAUDE.md): a statement concluding in BF₀ carries the LE reading alongside.

\section{Self-decomposability and the characterization theorem}\label{sec:characterization}

% [transcription pending]
